% ---- pg-269 （书页 257） ----

\begin{mvfigure}
\begin{tikzpicture}[x=7.5cm,y=1.45cm,>=Stealth,font=\zihao{6}]
  \draw[->] (3.87,0.78) -- (4.86,0.78);
  \draw[->] (3.87,0.78) -- (3.87,4.10);
  \foreach \xx/\lab in {3.9/3.9,4.0/4.0,4.1/4.1,4.2/4.2,4.3/4.3,4.4/4.4,4.5/4.5,4.6/4.6,4.7/4.7,4.8/4.8}{%
    \draw (\xx,0.78) -- (\xx,0.72);
    \node[below] at (\xx,0.72) {\lab};}
  \foreach \y in {1,2,3,4}{%
    \draw (3.87,\y) -- (3.82,\y);
    \node[left] at (3.82,\y) {\y};}
  \node[below] at (4.36,0.34) {$\ln(\text{SVL})$};
  \node[rotate=90] at (3.72,2.44) {$\ln(\text{重量})$};
  % S 组：空心圆
  \foreach \p in {(3.95,1.12),(4.02,1.45),(4.05,1.47),(4.08,1.50),(4.12,1.60),
    (4.15,1.68),(4.17,1.76),(4.20,1.92),(4.22,1.99),(4.25,2.06),
    (4.27,2.14),(4.30,2.24),(4.32,2.31),(4.35,2.39),(4.37,2.47),
    (4.40,2.54),(4.42,2.61),(4.45,2.68),(4.47,2.76),(4.50,2.86),
    (4.52,2.98),(4.55,3.10),(4.57,3.22),(4.60,3.32),(4.62,3.42),
    (4.65,3.52),(4.67,3.58),(4.68,3.45),(4.70,3.62),(4.72,3.55)}{%
    \draw[fill=white] \p circle (1.5pt);}
  % C 组：实心圆
  \foreach \p in {(4.03,0.78),(4.22,1.52),(4.25,1.55),(4.28,1.62),(4.30,1.70),
    (4.33,1.72),(4.35,1.80),(4.38,1.86),(4.42,1.98),(4.45,2.05),
    (4.48,2.15),(4.50,2.42),(4.53,2.52),(4.55,2.58),(4.57,2.62),
    (4.58,2.70),(4.60,2.62),(4.60,2.75),(4.62,2.80),(4.65,2.85)}{%
    \fill \p circle (1.5pt);}
  \draw (4.63,3.50) rectangle (4.87,3.98);
  \fill (4.67,3.80) circle (1.5pt);
  \node[anchor=west] at (4.71,3.80) {C};
  \draw[fill=white] (4.67,3.62) circle (1.5pt);
  \node[anchor=west] at (4.71,3.62) {S};
\end{tikzpicture}
\phantomsection\label{fig:6.7}\mvcap{图 6.7\quad 表 6.7 中蜥蜴数据的 $\ln(\text{重量})$ 对 $\ln(\text{SVL})$ 的散点图}
\end{mvfigure}

在对数据取了自然对数之后，汇总统计量如下：
\[
C:n_1=20,\quad \bar{\mathbf{x}}_1=\begin{bmatrix}2.240\\4.394\end{bmatrix},\quad
\mathbf{S}_1=\begin{bmatrix}0.353\,05&0.094\,17\\0.094\,17&0.025\,95\end{bmatrix}
\]
\[
S:n_2=40,\quad \bar{\mathbf{x}}_2=\begin{bmatrix}2.368\\4.308\end{bmatrix},\quad
\mathbf{S}_2=\begin{bmatrix}0.506\,84&0.145\,39\\0.145\,39&0.042\,55\end{bmatrix}
\]

\hyperref[fig:6.7]{图 6.7} 是取过自然对数后的重量和鼻子到排泄口长度（SVL）的散点图．$\ln(\text{重量})$ 的均值差和 $\ln(\text{SVL})$ 的均值差的大样本单－ $95\%$ 置信区间都包括零．
\[
\ln(\text{重量}):\ \mu_{11}-\mu_{21}:\ (-0.476,0.220)
\]
\[
\ln(\text{SVL}):\ \mu_{12}-\mu_{22}:\ (-0.011,0.183)
\]
对应的用来检验无差别个体均值的单变量学生氏 $t$ 检验统计量的 $p$ 值分别为 0.46 和 0.08．显然，通过单变量的透视，对于两类蜥蜴我们都无法检测到重量均值的差别和 SVL 均值的差别．

但是，和散点\hyperref[fig:6.7]{图 6.7} 一致的是，二元检验分析明显地说明两类蜥蜴在长度上有差别．利用结果 6.4（参见例 6.5），$T^2$ 统计量近似于 $\chi_2^2$ 分布．例如，$T^2=225.4$ 且 $p$ 值小于 $0.000\,1$．所以在这种情况下，多元的方法是必要的．\hfill$\blacksquare$

例 6.16 和例 6.17 阐明了多元检验与相应一元检验之间的相对功效，用例 6.1 的数据同样有这种情形出现．

当从多个总体抽取随机样本时（回顾 6.4 节的单因子 MANOVA），多元检验是以下列矩阵为依据：
\[
\mathbf{W}=\sum_{l=1}^{g}\sum_{j=1}^{n_l}(\mathbf{x}_{lj}-\bar{\mathbf{x}}_l)(\mathbf{x}_{lj}-\bar{\mathbf{x}}_l)'\ \text{与}\
\mathbf{B}=\sum_{l=1}^{g}n_l(\bar{\mathbf{x}}_l-\bar{\mathbf{x}})(\bar{\mathbf{x}}_l-\bar{\mathbf{x}})'
\]

在整个这一章中，我们使用的检验统计量为
\[
\text{威尔克斯 } \Lambda \text{ 统计量}\quad \Lambda^*=\frac{|\mathbf{W}|}{|\mathbf{B}+\mathbf{W}|}
\]

它等价于似然比检验．但在许多统计软件包的输出中，还经常会遇到下面三个多元检验统计量：

\noindent 劳雷-霍特林迹统计量 $=\operatorname{tr}[\mathbf{B}\mathbf{W}^{-1}]$

\noindent 皮莱迹统计量 $=\operatorname{tr}[\mathbf{B}(\mathbf{B}+\mathbf{W})^{-1}]$

\noindent 罗伊最大根统计量 $=\mathbf{W}(\mathbf{B}+\mathbf{W})^{-1}$ 的最大特征根

当样本容量非常大时，所有这四种检验方法几乎是等价的．对中等大小的样本容量，所有比较均以通过模拟来研究的少数最必要的案例为依据．从此刻为止所报导的那些模拟中可知，前三种检验具有相似的功效，而罗伊最大根检验却与它们有差别，仅当存在一个非零特征
